%% The first command in your LaTeX source must be the \documentclass command.
%%
%% Options:
%% twocolumn : Two column layout.
%% hf: enable header and footer.
\documentclass[
% twocolumn,
% hf,
]{ceurart}

%%
%% One can fix some overfulls
\sloppy

\usepackage{my-main-style}

%%
%% Minted listings support
%% Need pygment <http://pygments.org/> <http://pypi.python.org/pypi/Pygments>
\usepackage{listings}

%% auto break lines
\lstset{breaklines=true}

%%
%% end of the preamble, start of the body of the document source.
\begin{document}

%%
%% Rights management information.
%% CC-BY is default license.
\copyrightyear{2025}
\copyrightclause{Copyright for this paper by its authors.
  Use permitted under Creative Commons License Attribution 4.0
  International (CC BY 4.0).}

%%
%% This command is for the conference information
\conference{HyperAgents'25: 2nd Workshop on Hypermedia Multi-Agent Systems,
  October 25-30, 2025, Bologna, Italy}

%%
%% The "title" command
\title{The Gap Between BDI Agents and Semantic Hypermedia and What We Can Do About It}

% \tnotemark[1]
% \tnotetext[1]{You can use this document as the template for preparing your
%   publication. We recommend using the latest version of the ceurart style.}

%%
%% The "author" command and its associated commands are used to define
%% the authors and their affiliations.
\author[1]{Samuele Burattini}[%
orcid=0009-0009-4853-7783,
email=samuele.burattini@unibo.it,
]
\cormark[1]
\fnmark[1]

\author[1]{Martina Baiardi}[%
orcid=0009-0001-0799-9166,
email=m.baiardi@unibo.it,
]
\fnmark[1]

\author[1]{Giovanni Ciatto}[%
orcid=0000-0002-1841-8996,
email=giovanni.ciatto@unibo.it,
]


\author[1]{Danilo Pianini}[%
orcid=0000-0002-8392-5409,
email=danilo.pianini@unibo.it,
]



\address[1]{\disi, \unibo}


%% Footnotes
\cortext[1]{Corresponding author.}
\fntext[1]{These authors contributed equally.}

%%
%% The abstract is a short summary of the work to be presented in the
%% article.
\begin{abstract}
  Traditional BDI agents,
  rooted in logic programming,
  remain challenging to integrate with the hypermedia nature of open Web environments
  that rely on Semantic Web technologies like RDF and OWL,
  This paper examines the gap between these paradigms and surveys existing integration efforts on a conceptual and technical level 
  and argues that existing tools offer limited ergonomic support for developers.
  Our proposal for a deeper integration relies on a generalized BDI engine to enable the development of BDI agents that can directly reason and operate on semantic hypermedia resources.
  We derive ideal requirements and show an abstract architecture that can be tailored to support different types of beliefs and reasoning mechanisms.
\end{abstract}

%%
%% Keywords. The author(s) should pick words that accurately describe
%% the work being presented. Separate the keywords with commas.
\begin{keywords}
  BDI  \sep
  Hypermedia \sep
  Agent-Oriented Programming
\end{keywords}

%%
%% This command processes the author and affiliation and title
%% information and builds the first part of the formatted document.
\maketitle

%==========================================================
\section{Introduction}
%==========================================================

Integrating intelligent agents with the Web is a long-standing goal
in both the \ac{MAS}~\cite{DBLP:conf/edoc/ShafiqDF06}
and Semantic Web~\cite{lassila2001semantic} communities, but,
despite such common interest,
the famous question
``Where are all the intelligent agents?''
remains unanswered to this day~\cite{hendlerb2007expert}:
the Web is still dominated by traditional clients, servers,
and (most recently) bots and scrapers%
---which can hardly be considered full agents.
%
However,
as recent developments in Web engineering have led to
the pervasive use of the \ac{REST} architectural style~\cite{DBLP:journals/toit/FieldingT02},
and the growing adoption of Web interoperability standards % even outside the traditional Web domain
-- for instance the \ac{WoT}~\cite{wotarch} --,
the integration of \ac{MAS} and the Web gained new momentum.

Along this line,
the \ac{W3C} Autonomous Agents on the Web community group\footnote{\url{https://www.w3.org/community/webagents/}}
has been founded to connect researchers and practitioners
and a new perspective has been proposed: \ac{hMAS}~\cite{DBLP:conf/atal/CiorteaMGBRZ19}
envision the (Semantic) Web as an open environment in which
agents can interact by consuming and manipulating hypermedia resources.

Even more recently,
with the advent of \ac{LLM}-based \emph{Agentic AI}~\cite{acharya2025access}
the creation of \ac{LLM}-driven agents that interact with Web APIs~\cite{10.5555/3692070.3692540}
-- or directly with Websites~\cite{10.5555/3692070.3694608} --
has gained significant attention.
%
However,
despite being promising,
these approaches
do not make structured knowledge representation obsolete~\cite{pan2024tkde},
and often sacrifice the controllability of traditional agent-programming paradigms
%to sell the (hollow?) promise of circumventing the need for structured knowledge representation and reasoning~\cite{pan2024tkde}%
at the price of relaxed guarantees in terms of correctness, and completeness.
% provided by Semantic Web technologies such as the \ac{RDF}~\cite{RDF_Concepts_W3C:14} and \ac{OWL}~\cite{OWL_Syntax_W3C:12}.

Arguably,
Web-integrated agents developed with cognitive architectures,
such as \ac{BDI} agents,
would retain controllability and predictability:
they could directly exploit hypermedia to discover new resources and actions,
and use Semantic Web technologies like \ac{RDF}~\cite{RDF_Concepts_W3C:14} and \ac{OWL}~\cite{OWL_Syntax_W3C:12}
to act intelligently without compromising on quality.
%
However,
we believe
the integration of \ac{BDI} agents with the Web
is hindered by a conceptual and technological gap,
which limits developers in creating agents that need to interact with (Semantic) hypermedia.
%
More precisely:
\begin{enumerate*}[label=\textbf{(G\arabic*)}]
  \item\label{gap:logic}
  \ac{BDI} agents are historically tied to \ac{LP}
  %-- e.g., with the popular \agentspeak{} semantic~\cite{DBLP:conf/maamaw/Rao96} --
  whereas the Semantic Web is grounded on \ac{DL},
  and although the two paradigms are theoretically compatible,
  their practical implementations interoperate poorly,
  increasing integration cost significantly;
  % The different technological choices used to concretely implement them add cognitive load to developers working on the integration of the two worlds (e.g., converting \ac{OWL} and \ac{RDF} triples in Prolog-like formulas);

  \item\label{gap:open-world}
  \ac{BDI} agents are typically based on the \ac{PRS} architecture~\cite{georgeff1986pieee}
  which relies on pre-defined plans and does not natively support discovering new actions at runtime,
  a necessary feature to interact proficiently with hypermedia environments.
\end{enumerate*}

In this paper,
we reflect on these gaps,
their technical implications, % on the overall development experience of creating \ac{BDI} agents that interact with Web resources
explore how this gap is currently being addressed in the \ac{hMAS} community and
propose how to mitigate it by abstracting the design of agent-oriented programming languages and frameworks.
% %
% We start by providing a brief overview of the \ac{BDI} model,
% the hypermedia nature of the (Semantic) Web,
% and
% the fundamental principles of \ac{hMAS} in \Cref{sec:background}.
% %
% We then analyze the gap between \ac{BDI} agents and hypermedia,
% and identify ideal integration requirements in \Cref{sec:gap-analysis}
% %
% Then,
% in \Cref{sec:related-works},
% we test those requirements against prior contributions
% which already attempted to bridge the gap between \ac{BDI} agents and semantic hypermedia.
% %
% Finally,
% in \Cref{sec:generalized-bdi-engine},
% We present our research direction
% aimed at %tightening the integration of \ac{BDI} agents and hypermedia
% reducing the abstraction gap for developers
% by introducing a generalized \ac{BDI} engine
% that can directly leverage Semantic Web technologies for belief representation and deliberation;
% and we discuss open challenges and opportunities of our approach.

%==========================================================
\section{Background}
\label{sec:background}
%==========================================================

\subsection{BDI Agents}

The \ac{BDI} model~\cite{DBLP:conf/atal/GeorgeffPPTW98}
is a well-established cognitive architecture for the development of intelligent agents.
%
% It is based on the idea of representing an agent's mental state in terms of three key components:
% \begin{inlinelist}
%   \item \emph{beliefs} are what the agent thinks to be true,
%   \item \emph{desires} are the potential goals or objectives an agent may pursue,
%   \item \emph{intentions} are the agent's commitments to achieve specific desires.
% \end{inlinelist}
\ac{BDI} agents follow a continuous cycle of
\begin{inlinelist}
  \item \emph{perception} of external stimuli and gather new \emph{beliefs},
  \item \emph{deliberation} to update beliefs, goals and choose \emph{intentions} to pursue selecting \emph{plan} to achieve goals,
  and
  \item \emph{action} advancing one intention by executing a step of the associated plan.
\end{inlinelist}
Practically, most \ac{BDI} agents implementations follow the \ac{PRS} architecture~\cite{georgeff1986pieee},
which is based on the ability for agents to select plans from a predefined set
and execute them to achieve their goals.

To declaratively define the agent's behavior the \ac{AOP} community has developed several \ac{BDI} programming languages and frameworks~\cite{DBLP:conf/woa/MascardiDA05}.
%
Arguably, the most popular language is \agentspeak{(L)}~\cite{DBLP:conf/maamaw/Rao96} which
uses a formal syntax and semantics derived from \ac{LP} for representing beliefs, goals and actions using logic predicates.
%This allows to exploit logic unification to query the agent's \emph{belief base} and match plans to achieve goals.
%The \agentspeak{} abstract language has been implemented in the \jason{}~\cite{bordini2007programming} agent programming platform, providing an interpreter and a Java-based runtime environment for developing \ac{BDI} \ac{MAS}.

% The logic-based nature of \agentspeak{} has characterized the development of \ac{BDI} agents due to its expressiveness and flexibility in supporting the basic mechanisms of the \ac{BDI} reasoning cycle.
%
Despite \ac{BDI} programming languages and frameworks have also been implemented without such a strong reliance on \ac{LP} mechanisms (e.g., \cite{map2005sp,kampik2019emas}), in the following when we refer to \ac{BDI} agents we will implicitly refer to \agentspeak{(L)}-like agents.

\subsection{The Web, Semantics and Hypermedia}

The World Wide Web was proposed by Tim Berners-Lee as a global, open information space
based on the ability to link resources through hypermedia documents~\cite{DBLP:journals/cn/Berners-LeeCG92}.
%
The Web is built with the \ac{REST} architectural style~\cite{DBLP:journals/toit/FieldingT02}
which supports independently evolving components through a set of principles.
%
% The \ac{REST} architectural style~\cite{DBLP:journals/toit/FieldingT02}
% adopted by the Web allows to have independently evolving components
% that can interact through a uniform interface thanks to:
% \begin{inlinelist}
%   \item \acp{URI} to identify resources,
%   \item the exchange of representations of resources using shared data formats (e.g., JSON, XML),
%   \item self-descriptive messages following the semantics of the HTTP protocol, and
%   \item the use of \ac{HATEOAS} guiding clients in their next possible actions.
% \end{inlinelist}

The Semantic Web~\cite{BernersLee2001} complements the \emph{self-descriptive} principle of \ac{REST}
by encoding semantics in the representation through the general model of \ac{RDF}~\cite{RDF_Concepts_W3C:14}
based on machine-readable \emph{triples}
that state facts using a common vocabulary defined within ontologies.
%
{OWL}~\cite{OWL_Syntax_W3C:12} allows the encoding of such ontologies following the principles of \ac{DL} to define concepts and relationships between them.
%
%Ontologies can be used to encode structured knowledge about a domain of interest, creating consensus on the meaning of concepts and relationships between them, hence improving interoperability between different systems on the open Web.

Another fundamental principle is \ac{HATEOAS} which defines hypermedia as ``the simultaneous presentation of information and controls such that the information becomes the affordance through which the user obtains choices and selects actions''~\cite{DBLP:journals/toit/FieldingT02}.
% %
% Affordances
% -- i.e., action possibilities --
% are central to the inner workings of the Web,
% allowing clients to discover new resources and actions by consuming hypermedia.
% This mechanism drives the state transitions of the distributed client-server system
% ensuring scalability and independence of individual components.
In this way, the Web can be seen as an open environment
where clients discover new resources and action possibilities by consuming hypermedia.
%
When such hypermedia is annotated with explicit semantics it should be possible for clients to better understand the meaning of the affordances, avoiding the need for hard-coded knowledge.
%
%This is the mechanism adopted for instance in the \ac{WoT}~\cite{wotarch} where \emph{\ac{TD}} are used as semantic hypermedia representations of \ac{IoT} devices,
%declaring the possible interaction affordances exposed by the devices (e.g., reading a sensor value, turning on a light, etc.).


\subsection{Hypermedia Multi-Agent Systems}

The idea of \ac{hMAS}~\cite{DBLP:conf/atal/CiorteaMGBRZ19} seeks a deeper integration of \ac{MAS} and the Web through the adoption of hypermedia and \ac{REST} principles.
%
The main idea behind such integration is to overcome the limitations of previous approaches to \ac{MAS} and Web integration, by considering the Web as an \emph{environment} for agents to operate in~\cite{ciortea2018emas}.
%
Following the idea of the environment as a first-class abstraction of \ac{MAS}~\cite{weyns2007aamas},
\ac{hMAS} revisits the Agents and Artifacts meta-model~\cite{ricci2011aamas} as a way to modularize and program the environment.
%
In \ac{hMAS}, artifacts are represented as hypermedia resources, exposing affordances that agents can discover through the uniform interface of the Web.

Most of the effort in the \ac{hMAS} community has been focused on shaping the environment dimension of the \ac{MAS}, as ideally heterogeneous agents should be able to operate in the same environment.
%
However, the agent dimension of \ac{hMAS} must also be investigated to improve how agents can effectively exploit the hypermedia environment (e.g., for action discovery at runtime~\cite{vachtsevanou2024atal}).
%
In this paper, we focus on such agent dimension of \ac{hMAS} and specifically on \ac{BDI} agents and how to support deeper integration with semantic hypermedia.

%==========================================================
\section{Integrating \acs{BDI} Agents and Semantic Hypermedia: a Gap Analysis}
\label{sec:gap-analysis}
%==========================================================

% From the brief background presented in \Cref{sec:background},
% it should be evident that
% despite both \ac{BDI} agents and the Semantic Web being rooted in logic-based knowledge representation,
% there are significant differences between the two worlds.
% %Specifically, Prolog-like logic typically used to practically implement \ac{BDI} agents and \ac{DL} used in the Semantic Web are different subsets of \ac{FOL} with consequently different expressiveness and semantics
% Specifically, \ac{DL} excel in representing relationships between concepts, while \ac{BDI} agents use \ac{LP} to store facts about the external environment as belief and declaratively express goals.
% The different technologies used in the two different paradigms makes them difficult to integrate \ref{gap:logic}.
% %
% Furthermore, \ac{BDI} agents are typically designed with a closed-world perspective whereas the Semantic Web is an open information space \ref{gap:open-world}.
% %
% When it comes to practical applications this creates an abstraction and technological gap
% which can be difficult to bridge for developers.
% %
% In the following, we describe in more detail such differences.

Here we delve into the details of the gaps that we identified between \ac{BDI} agents and semantic hypermedia,
namely,
\ref{gap:logic} the practical differences for knowledge representation and manipulation
between \ac{LP} and Semantic Web technologies,
and
\ref{gap:open-world} the limitations of \ac{BDI} agents with \emph{open} hypermedia environments,
where actions are not predefined but discovered at runtime.

\paragraph{Syntactical Differences Between \acs{LP} and \acs{RDF}.}

On a syntactical level,
\ac{BDI} agents commonly represent beliefs as definite clauses~\cite{MAKOWSKY1987266}%,
% and goals as facts
%which are either atomic \emph{terms} (e.g., \texttt{alice}) or a \emph{structure} composed of a \emph{functor} followed by an arbitrary number of \emph{arguments} within parentheses
(e.g. \texttt{person(alice)}).
%
In the Semantic Web,
knowledge is represented in the form of \ac{RDF} \emph{triples} which connect a \emph{subject} to an \emph{object} through a \emph{predicate} (e.g. \texttt{:alice rdf:type foaf:Person}).
%
Although \ac{RDF} triples can be represented as definite clauses~\cite{DBLP:conf/www/GrosofHVD03}
(while the opposite is not true),
in practice,
semantic knowledge is available on the Web in \ac{RDF},
and accessing that from \ac{BDI} agents requires a translation step,
which has many degrees of freedom
(e.g., \texttt{type(alice, person)} and \texttt{triple(alice, type, person)} are reasonable mappings of the \ac{RDF} triple above),
and it is up to \ac{BDI} developers.
%
% Developers wishing to integrate \ac{BDI} agents with the Semantic Web hence face the challenges of choosing a suitable -- and more importantly consistent -- mapping which we believe adds a significant cognitive load during the development process.
In other words,
implementing the translation is additional work which, if handled manually in an ad-hoc manner can be source of inconsistencies.

\paragraph{Ontological Reasoning.}

\ac{RDF} triplets are commonly paired with \ac{OWL} ontologies (e.g., \texttt{foaf}\footnote{\url{http://xmlns.com/foaf/spec/}} in the example above),
providing \emph{axioms} that define the meaning of the triples's predicates
(e.g., \texttt{rdf:type} denotes the \emph{instance of} relation between subjects and classes),
as well as constraints, and relationships between concepts.
%
This is a way to standardize knowledge representation across hypermedia.
%
Logic programs, instead, are not really meant to be scattered over the Web,
nor being accessed remotely.

Again,
there is a hidden trap for \ac{BDI} developers willing to create agents that can perform ontological reasoning:
% This means that, in order to properly reason about the knowledge represented in \ac{RDF} and \ac{OWL},
% developers may need to translate the whole \ac{RDF}, \ac{OWL} and ontology semantics into their own belief base, adding a significant overhead and a potential source of inconsistencies.
either they translate the \ac{OWL} ontology into definite clauses,
or they wrap some \ac{OWL} reasoner into the agent interpreter and reasoning cycle.
%
Despite being automatable to some extent (cf. ~\cite{samuel2008tplp}),
both activities require additional engineering effort,
and should rather be addressed with direct support from the \ac{BDI} framework.

% Although \ac{OWL} ontologies can be converted into logic formulas through an automatic process
% it can be cumbersome for \ac{BDI} developers to exploit them when writing the agent program
% and similarly hard for Semantic Web experts to understand how to properly interact with the translated knowledge base.

\paragraph{Querying and Inference.}

\ac{BDI} commonly relies on \ac{SLD} resolution~\cite{emden1976jacm} to query the belief base,
when this is implemented as a logic program.
%
Furthermore,
they rely on logic unification~\cite{10.1145/357162.357169} to match plans and goals,
making the intertwining with \ac{LP} even deeper.
%
% This basic mechanism allows to easily verify conditions based on the current beliefs, retrieve information and also perform inference using dedicated rules developers can encode to represent the problem domain.
% %
% Unification relies on finding a substitution for the variables in a logic formula that makes it true with the facts within a logic theory. When multiple substitutions are possible, the unification process usually returns the first available one, as it is typically implemented as a depth-first search algorithm with backtracking.
%
Conversely,
the Semantic Web relies on \ac{SPARQL}~\cite{sparqlprotocol} to query \ac{RDF} data.
%
\ac{SPARQL} allows matching graph patterns and retrieve all the values of variables in the pattern from the matched triples,
providing an intuitive way to query the knowledge base.

Despite both \ac{LP} and \ac{SPARQL} allow expressing declarative queries for a knowledge base,
the latter may be more practical and efficient for \ac{RDF} data,
especially when the query spans over multiple triples.
%
Consider for instance a query searching for the unknown relations \texttt{?r} between \texttt{alice} and \texttt{bob},
and matching all the triples where the subject is \texttt{carl} and
the predicate is \texttt{?r} to retrieve all the objects \texttt{?x}.
%
A query on this kind would be straightforward to write in \acs{SPARQL},
while requiring second-order variables in \ac{LP}%
---a feature which is not commonly supported by \ac{LP} or \ac{BDI} frameworks,
and which is commonly worked around via \emph{meta-programming},
resulting in longer and cumbersome \ac{LP} code.

%
As for inference,
given the expressiveness limitations of \ac{OWL},
\ac{SWRL}~\cite{horrocks2004swrl} can be used to define inference rules.
%
% An important difference is that while logic unification automatically performs inference, \acs{SPARQL} searches only on the set of triples currently in a knowledge base. If inference is required, the developer must explicitly \emph{materialize} the triples inferred by a reasoner before querying.
Here, a similar argument like the one above concerning the wrapping of \ac{SWRL} reasoners applies.

\paragraph{Open vs. Close Environments.}

\ac{BDI} agents are equipped with a predefined set of plans, composed of sequences of actions that the agent knows how to execute.
%
Additionally, \ac{BDI} agents usually run in a closed environment,
which is designed by the developer to provide perceptions to the agents and
support the execution of the actions used in the plans~\cite{weyns2007aamas}.
%
The development process of \ac{BDI} agents is hence based on the fact that a developer anticipates
the relevant possible settings an agent may encounter and provides plans to handle them.
%
This makes it difficult for agents to adapt to unexpected settings or to flexibly incorporate new actions that become available while exploring the environment.
%
Relating this to the open hypermedia nature of the Web, makes the effective integration of \ac{BDI} agents challenging.
%
Following the \ac{HATEOAS} principle, agents may discover new affordances
-- i.e., action possibilities --
while navigating through the hypermedia environment.
%
The inability to exploit such new affordances may limit the agent in adapting to the environment and effectively interact with it.

\subsection{Integration Requirements}
\label{sec:integrating-bdi-hypermedia}

Guided by our gap analysis,
we identify a set of requirements that we believe may improve practical development of \ac{BDI} agents in Web environments.
%
We hence propose that a \ac{BDI} agent programming framework for \ac{hMAS} should support the following features:
%
\begin{enumerate*}[label=\textbf{(R\arabic*)}]
  \item direct manipulation of \ac{RDF} and \ac{OWL} triples in the agent's belief base;
  \label{req:direct}

  \item ontological reasoning and inference for deliberation (e.g., plan selection and execution);
  \label{req:reasoning}

  \item support for querying the belief base via \acs{SPARQL};
  \label{req:query}

  \item ability to exploit affordances discovered in the environment to dynamically adapt the agents' plans.
  \label{req:actions}
\end{enumerate*}


%==========================================================
\section{Related Works}
\label{sec:related-works}
%==========================================================

% \sbnote{ho tolto i paragrafi ma non ho integrato, passo la palla a gio prima di conflicts}


%\sbnote{Related works, JASDL (Jason+Ontological reasoning), Yggdrasil framework, Hypermedea (Saint-Etienne Cartago Artifacts for Hypermedia), others.... }

Here we summarize some related works demonstrating how,
despite the many extensions of \agentspeak{} towards \ac{DL},
the direct exploitation of Semantic Web technologies
-- namely, \ref{req:direct} and \ref{req:query} --
in \ac{BDI} \emph{programming} frameworks is still underexplored.

% \paragraph{Semantic extensions of \agentspeak{}.}

In \cite{DBLP:conf/dalt/MoreiraVBH05} authors propose \agentspeakdl{},
an extension of \agentspeak{} to support the $\mathcal{ALC}$ description logic~\cite{BAADER2008135}.
%
% In \agentspeakdl{},
% semantic rules for the steps in the reasoning cycle are the same of predicate-logic \agentspeak{},
% but the beliefs representation is extended with \ac{DL} objects and relationships to express \ac{RDF} triples.
% %
% Their formal semantics description highlights
% that introducing ontological reasoning
% affects mainly three operations of the whole \agentspeak{} lifecycle, namely:
% \begin{inlinelist}
%   \item the Plan research,
%   \item the query on the Belief Base, and
%   \item the Belief Update.
% \end{inlinelist}
There,
agents' belief bases consist of an $\mathcal{ALC}$ ontology
(concepts, roles, and individuals)
and any operation involving access to the belief base
(e.g., plan selection, belief update, and query)
imply either reasoning about or updating such ontology.
%
So,
\agentspeakdl{} provides the theoretical foundations to support our integration requirements from \Cref{sec:integrating-bdi-hypermedia}.

\agentspeakdl{} is then extended by CooL-\agentspeak{}~\cite{DBLP:journals/wias/MascardiABBR14},
where authors explore inter-agent exchange of procedural knowledge~\cite{DBLP:conf/dalt/AnconaM03},
possibly expressed by different ontologies.
%
% To do so,
% CooL-\agentspeak{} integrates Coo-BDI~\cite{DBLP:conf/dalt/AnconaM03} and \agentspeakdl{}
% features.

The authors of \cite{DBLP:conf/iat/HalacED11} provide a different theoretical integration of \agentspeak{} with \ac{DL}
to account for \emph{goal states}
(new, suspended, active, succeeded, failed)
in the agent's lifecycle,
in such a way that state transitions are tied to conditions expressed in \ac{DL}.
% claim that \agentspeakdl{} does not provide a complete solution for the integration of \ac{BDI} agents and Semantic Web technologies,
% mainly because goals are only represented as events for procedural plans.
%
% Therefore,
% they propose an extension of the \agentspeak{} semantics
% where the agents lifecycle is enriched with goal states,
% where beliefs are logic predicates adapted to represent \ac{DL} objects and relationships,
% like in \agentspeakdl{}.
% \gcsidenote{Riprendi da qui}

% \paragraph{Integrating \ac{DL} and \ac{BDI}.}

\agentspeakdl{} is also the foundation for JASDL~\cite{DBLP:conf/dalt/KlapiscakB08},
i.e., an implementation of \agentspeakdl{} in \jason{}~\cite{bordini2007programming}.
%
The key extension in JASDL are ``semantically-enriched literals'', 
i.e. $\mathcal{ALC}$ statements, 
expressed as logic facts,
with ordinary Jason syntax for beliefs
(see example in \Cref{lst:jasdl-example}).
%
% The examples in \Cref{lst:jasdl-example} illustrate how \ac{RDF} triples are represented as \ac{DL} in the belief base of the agent.
%
JASDL uses \jason{}'s standard belief-base management mechanisms,
while wrapping and adapting a semantic reasoner to let agents perform ontological reasoning,
thus addressing \ref{req:reasoning}.
%
This approach gives developers the full potential of \jason{} plus \ac{DL}, 
at the (non-neglible) price of translating \ac{RDF} triples into JASDL-compatible logic predicates.
%
\lstinputlisting[float,
  lastline=3,
  language=Java,
  caption={On the left, \ac{RDF} triples expressed as predicates in JASDL, on the right their meaning inferred by a \ac{OWL} reasoner. The example is taken from \cite{DBLP:conf/dalt/KlapiscakB08}.},
  label={lst:jasdl-example}
]{listings/jasdl.txt}

%
Other approaches~\cite{DBLP:journals/wias/MascardiABBR14,DBLP:conf/www/CharpenayZLB22, DBLP:conf/semweb/ONeillBBC21, DBLP:conf/emas/CiorteaBR18}, rely on encapsulating functionalities and services in external components that agents can share and exploit to support their activities.
% where artifacts are a set of passive resources and tools encapsulating functionalities and services
% that agents can share and exploit to support their activities.
% using artifacts to mediate the interaction among agents and the semantic hypermedia environment.
%
For instance
CooL-\agentspeak{}~\cite{DBLP:journals/wias/MascardiABBR14} relies on \cartago{}~\cite{DBLP:books/sp/09/RicciPVO09},
to define and implement the \emph{Ontology Artifact},
providing agents with ontological reasoning capabilities%
---encoding \ac{DL} statements as beliefs similarly to JASDL.

Similarly, 
Hypermedea~\cite{DBLP:conf/www/CharpenayZLB22} extends \jacamo{}~\cite{BOISSIER2013747}
-- again, via artifacts and constraints on the beliefs-base syntax --
to support agents interacting with semantic hypermedia and, in particular, the \ac{WoT}.
%
More precisely,
Hypermedea \emph{artifacts} produce agent beliefs as the ones in \Cref{lst:hypermedea-example}
through observable properties
---i.e., \ac{RDF} triples expressed as logic predicates,
in Jason's syntax.
%
Furthermore, 
other \emph{artifacts} support the manipulation of such \ac{RDF} triples,
materialization of ontological inferences,
as well as the interaction with \acs{WoT}'s \emph{things},
and the synthesis of plans via a PDDL planner
---hence addressing \ref{req:actions}.
% \begin{inlinelist}
%   \item \textit{Linked Data Artifact} to interact with \ac{RDF} triples,
%   \item \textit{Ontology Artifact} which bridge the gap between \ac{RDF} models and predicate representations,
%   \item \textit{Thing Artifact} which proxies the communication with Things in the environment (i.e. each Thing its own Thing Artifact to interact with it), and
%   \item \textit{Planner Artifact} which synthetize plans given a \textit{domain description} and a \textit{Problem Description}, using PDDL definitions and FF-X planner.
% \end{inlinelist}
%
\lstinputlisting[float,
  % lastline=1,
  language=Java,
  caption={The example shows how \ac{RDF} triples are represented in Hypermedea agents belief base. The example is taken from Hypermedea \ac{API} reference (https://archive.is/vxULu).},
  captionpos=b,
  label={lst:hypermedea-example}
]{listings/hypermedea.txt}

A different take is proposed in \cite{DBLP:conf/semweb/ONeillBBC21}
where Astra~\cite{CollierRL15} agents are equipped with a set of \emph{modules} to allow \ac{BDI} agents to store knowledge and interact with a ``personal'' triple store separated from the belief base.
Similarly to other approaches,
the bridge to belief representation is by mapping triples directly to logic predicates with three arguments.

Finally, although not directly related to \ac{BDI} agents, the Yggdrasil framework~\cite{ciortea2018emas} provides a hypermedia layer for building \ac{hMAS} environments on top of \cartago{} artifacts
and has been used to implement \ac{hMAS} applications with the \jacamo{} framework.

% Finally,
% in line with the vision of the environment as a first-citizen abstraction for the interaction between hypermedia and agents~\cite{DBLP:conf/www/CharpenayZLB22},
% the authors of~\cite{DBLP:conf/emas/CiorteaBR18} propose a demonstrator application,
% % together with the engineering principles that support the integration using the environment.
% based on Yggdrasil~\cite{ciortea2018emas} for supporting the hypermedia abstraction,
% and on \jacamo{}~\cite{BOISSIER2013747} for implementing agents and artifacts.

\paragraph{Exploiting Discovered Actions at Runtime.}

To exploit new actions discovered at runtime \ref{req:actions} \ac{BDI} agents must have access to a \emph{planner} either within the agent -- see ~\cite{MeneguzziS15} for a general overview of planning in \ac{BDI} agents -- or through the environment---as proposed in Hypermedea~\cite{DBLP:conf/www/CharpenayZLB22}.
%
In the context of \ac{hMAS}, 
this requires that affordances in the environment 
formally describe the criteria for how and when to use them 
(e.g., their preconditions and effects),
and that agents can interpret such information to instruct the planner accordingly.
%
Recent works have explored this idea by using \emph{signifiers}~\cite{vachtsevanou2023atal} in hypermedia environments to better convey affordances to \ac{BDI} agents~\cite{vachtsevanou2024atal}.
%
A less structured approach is proposed in \cite{10.1145/3603163.3609077,schmid2024kgswc},
where agents (albeit not \ac{BDI}) are able to select actions in a hypermedia environment by using an \ac{LLM} oracle.

%==========================================================
\section{Levelling the Gap with a Generalised \ac{BDI} Engine}
\label{sec:generalized-bdi-engine}
%==========================================================

To bridge the gap between \ac{BDI} agents and semantic hypermedia,
we propose a \emph{generalized} \ac{BDI} engine,
addressing the requirements in \Cref{sec:integrating-bdi-hypermedia}.
%
This approach primarily tackles gap~\ref{gap:logic}, while leaving the exploration of gap~\ref{gap:open-world} -- benefitting from the modularity of the engine -- to future works.

The core idea is to create a \ac{BDI} interpreter encapsulating agents reasoning cycle,
while decoupling the specific technology used for knowledge representation and manipulation.
%
This would allow for tailoring the reasoning mechanisms to application-specific needs.
%
In this way,
beliefs
-- and, consequently, reasoning --
can flexibly support different formats depending on the domain at hand.
%
In this way,
agents can easily be adapted to work with different data formats,
ranging from \ac{RDF} triples, to JSON, YAML, and other Web standards.
%
Such a design fulfills the first requirement \ref{req:direct}.

\begin{wrapfigure}{R}{.35\linewidth}
    \centering
    \caption{Generalized Architecture.}
    \label{fig:generalized-bdi-architecture}
    \includegraphics[width=\linewidth]{figures/architecture.pdf}
\end{wrapfigure}

However,
separating the reasoning cycle of \ac{BDI} agents from knowledge representation is complex,
as the two are deeply intertwined~\cite{DBLP:conf/dalt/MoreiraVBH05}.
%
Consider for instance these fundamental operations:
%
\begin{inlinelist}
  \item \emph{plan selection}
  % is the process of selecting a plan to achieve a goal among the ones available,
  % which evaluation of applicability is strictly related on how beliefs are represented in the belief base;
  depending on conditions to be tested against the belief base;

  \item \emph{belief querying}
  % is the process of retrieving information from the belief base;
  as part of a plan, for the sake of decision-making;

  \item \emph{belief update}
  % is the process of updating the belief base with new information.
  for the sake of updating the agent's knowledge,
  as part of a course of action.
\end{inlinelist}
%
In \ac{LP}-based frameworks,
these would simply rely on logic unification and resolution,
whereas in a generalized \ac{BDI} engine,
each operation may rely on different matching strategies.
%
For instance,
for Semantic Hypermedia agent,
one may rely on \ac{RDF} and \ac{OWL} for beliefs representation,
\acs{SPARQL} for belief querying and update,
and \ac{SWRL} for plan selection and reasoning%
%
% In the latter case,
% agent developers would be able to plug-in convenient syntaxes and ad-hoc technologies,
% hence customising how the generic \ac{BDI} engine operates upon beliefs in all such cases,
---thus addressing \ref{req:reasoning} and \ref{req:query}.

% The integration of \ac{BDI} \ac{MAS} with heterogeneous syntax representations
% can be performed with a layered architecture of the framework.
% %
% The layered architecture proposed is composed of:
Such a generalized \ac{BDI} engine can be implemented with a layered architecture,
depicted in \Cref{fig:generalized-bdi-architecture}, and composed of three main layers.
%
Firstly,
\begin{enumerate*}[label={\arabic*.}]
  \item the \textbf{generalized \ac{BDI} reasoning layer},
  where the agent's deliberation process is implemented as a reasoning cycle,
  yet agnostic to how beliefs are represented and manipulated.
  %
  This is where,
  agents' intentions are maintained
  and updated as agents perceive,
  decide what to do,
  and finally act.
  %
  Any operation involving belief bases should rely on some \emph{abstract} \acs{API},
  to be implemented by the next layer.
  %
  \item The \textbf{concrete specification layer}:
  where the aforementioned \acs{API} is implemented to provide concrete operations on the belief base.
  %
  This layer acts as an adapter between the generalized \ac{BDI} reasoner
  and the specific knowledge representation technologies of choice.
  %
  Lastly,
  \item the \textbf{application layer}:
  is where actual goals, plans, and interactions of a \ac{MAS} are defined
  to tackle some specific use case.
\end{enumerate*}

% The main advantages of this isolation are: reuse of the same reasoning engine across multiple domains, and Better support for semantic web standards and open-world reasoning.

Requirement \ref{req:actions} although not directly addressed by the approach,
could be integrated with the techniques already being explored in the \ac{hMAS} community (cf. \Cref{sec:related-works}), benefitting from the possibility of choosing the most suitable knowledge representation technology
to facilitate their agents to discover new actions at runtime in the hypermedia environment.

%==========================================================
\section{Conclusion}
\label{sec:conclusion}
%==========================================================

In this paper, we present a gap analysis between \ac{BDI} agents and semantic hypermedia,
showing how the two worlds are difficult to integrate due to the different technologies and paradigms they rely on.
%
We identify a set of ideal requirements, finding that none of the surveyed related works fully satisfy them
and propose to implement them through a generalized \ac{BDI} engine that can be tailored to use \ac{RDF} and \ac{OWL} as knowledge representation technologies.

Future works will focus on implementing the architecture proposed in \Cref{sec:generalized-bdi-engine}
and explore how to support \ac{hMAS} developers with syntax extensions and tools to ease the development of \ac{BDI} agents in hypermedia environments within our \ac{BDI} programming framework \jakta{}~\cite{baiardi2024sncs}.


% \sbnote{Challenges and opportunities of the approach, e.g. ontological reasoning and inference, Belief consistency, "goal" definition and management, }

% \gcnote{è uno short position paper, va bene anche senza conclusioni, o se le conclusioni abbozzano semplicemente le sfide e le opportunità del framework proposto.}

%%
%% The acknowledgments section is defined using the "acknowledgments" environment
%% (and NOT an unnumbered section). This ensures the proper
%% identification of the section in the article metadata, and the
%% consistent spelling of the heading.
\begin{acknowledgments}
  This work has been partially supported by:
  %
  \begin{inlinelist}
      \item ``WOOD4.0 - Woodworking Machines for Industry 4.0'',
          Emilia-Romagna CUP E69J22007520009;
      \item ``FAIR--Future Artificial Intelligence Research'',
          Spoke 8 ``Pervasive AI''
          (PNRR, M4C2, Investimento 1.3, Partenariato Esteso PE00000013),
          funded by the EC under the NextGenerationEU programme;
      \item “ENGINES — ENGineering INtElligent Systems around intelligent agent technologies” project funded by the Italian MUR program ``PRIN 2022'' (G.A.\ 20229ZXBZM),
      and
      \item 2023 PhD scholarship
          (PNRR M4C2, Investimento 3.3 DM 352/2022),
          co-funded by the European Commission
          and AUSL della Romagna.
  \end{inlinelist}
\end{acknowledgments}

%% The declaration on generative AI comes in effect
%% in Janary 2025. See also
%% https://ceur-ws.org/GenAI/Policy.html
\section*{Declaration on Generative AI}

%   {\em Either:}\newline
%   The author(s) have not employed any Generative AI tools.
%   \newline

%  \noindent{\em Or (by using the activity taxonomy in ceur-ws.org/genai-tax.html):\newline}
%  During the preparation of this work, the author(s) used X-GPT-4 and Gramby in order to: Grammar and spelling check. Further, the author(s) used X-AI-IMG for figures 3 and 4 in order to: Generate images. After using these tool(s)/service(s), the author(s) reviewed and edited the content as needed and take(s) full responsibility for the publication’s content.
During the preparation of this work,
the authors used GitHub's Copilot in order to spell-check their writing, and make it more concise.
%
After using this tool,
the authors reviewed and edited the content as needed and take full responsibility for the content of the publication.
%%
%% Define the bibliography file to be used
\bibliography{bibliography}

%%
%% If your work has an appendix, this is the place to put it.
% \appendix

% \section{Online Resources}
% The sources for the ceur-art style are available via
% \begin{itemize}
% \item \href{https://github.com/yamadharma/ceurart}{GitHub},
% % \item \href{https://www.overleaf.com/project/5e76702c4acae70001d3bc87}{Overleaf},
% \item
%   \href{https://www.overleaf.com/latex/templates/template-for-submissions-to-ceur-workshop-proceedings-ceur-ws-dot-org/pkfscdkgkhcq}{Overleaf
%     template}.
% \end{itemize}

\end{document}

%%
%% End of file
